For out base neural network, we use a fully-connected network $N_{\theta}: (x^*, t^*) \rightarrow (\rho^*, u^*, p^*, T^*)$, with three hidden layers of $128$ neurons and $\tanh$ activations, and weights initialized with a Kaiming-normal scheme matched to the $\tanh$ non-linearity (zero-initialized biases).

Since dimensionless output properties are physically positive in nature, we apply a hard-constraint on each output channel by passing it through a shifted soft-plus function with a small constant $\epsilon$ added, to ensure that we get strictly positive $\rho^*, p^*, T^*$ and non-negative $u^*$. This forces them, instead of relying on soft loss terms to discourage behaviour that goes against the physics of our problem.

\begin{align*}
    \hat{q}_\text{out} &= softplus(\hat{q} + \epsilon), \quad q \in \{ \rho, p, T \} \\
    \hat{u}_\text{out} &= softplus(\hat{u} + \epsilon)
\end{align*}

\subsection{Base Residuals and Loss Functions}

From every batch of collocation points, the derivatives of the network output are obtained from passing them to a residual function:

\begin{align*}
    \mathcal{L}_\text{total} = \mathcal{L}_\text{PDE} + \mathcal{L}_\text{BC}
\end{align*}

where $\mathcal{L}_\text{PDE}$ stands for the property balances according to the governing equations themselves. However, boundary conditions are broadly categorized into a few different types. They are weighed heavily, since their penalties are evaluated on points drawn from the thin boundary layers of the domain (across a half-width $\delta = 10^{-5}$). Initially, we consider the specific given conditions:

\begin{enumerate}
    \item Inlet $(x^* = 0)$ - matching the respective dimensionless numbers for $p, \rho, T \space \text{and} \space u$.
    \item Outlet $(x^* = 1)$ - matching the back pressure, $p^*_b$.
    \item Initial condition $(t^* = 0)$ - matching the start-up conditions of the nozzle.
\end{enumerate}

Later, as detailed in section , we add on more conditions after deciphering the shock location (domain boundary):

\begin{enumerate}
    \item Across the interface, balance of mass, momentum and energy on either side of $\hat{x}_\text{shock}$.
    \item At the interface - properties may get completely discontinuous, inspite of having the balance maintained.
\end{enumerate}

Collocation points for building the residuals are periodically resampled per causal weighing (ref) and residual based adaptive refinement (ref).

\subsection{Stage I - Global Pre-training and Sub-network initialization}
\label{subsec:stage-1-global-pretrain-and-subnet-init}

Similar to [Addpinn ref], a single instance of the network mentioned in section , is trained over the entire domain $[x_\text{inlet}, x_\text{outlet}] \times [0, t_\text{steady}]$ using the overall boundary conditions. Reactive weighing is also performed for the boundary condition losses, to compensate for the lack of collocation points. A two-phase training is performed: an longer ADAM warm-up phase with gradient clipping to prevent gradient explosion and causal weighed point resampling, as well as a shorter LBFGS refinement phase with residual based adaptive refinement and a fresh resampling of points for non-finite loss occurrences.

After training, the network is evaluated at $t = t_\text{steady}$ throughout the grid. The derivatives, computed by automatic derivation, are used for the localization of shock - from both maxima and minima, the gradient is seen by walking backward until the slope drops/rises below a certain $\epsilon_s$. This yields a range of locations from which an initial guess of the shock location can be appropriated.

\subsection{Stage II - Domain decomposed training with boundary parameters}

The domain is then decomposed into a pre-shock $[x_\text{inlet}, \hat{x}_\text{s}]$ subdomain and a post-shock $[\hat{x}_\text{s}, x_\text{outlet}]$ subdomain, both being assigned with the same network architecture and warm-started with the weights obtained from the global, converged network.

Rather than a fixed interface, the location of the shock itself becomes a scalar parameter that can be learned and clamped within the nozzle length. The position of the discontinuity is free ro shift as the two sub-networks converge to the solution on both sides. Each sub-network is subject to residual losses in its own sub-domain: inlet and initial conditions for the pre-shock region and outlet and initial conditions for the post-shock region, besides the relation between properties across the sub-domain boundaries.

We follow the same two-phase training process as in stage I, but now, we simultaneously train both sub-networks with the shock location parameter. In addition to fresh resampling of points, we sample the sub-domain boundaries themselves as the shock location parameter shifts. This residual-based re-partitioning actually ensures sharper convergence at the cost of encountering race conditions.

\subsection{Interface coupling: Rankine-Hugoniot Jump Condition}

In order to reinforce the sub-networks around the shock location itself, a new set of losses are calculated on the points around $[\hat{x}_\text{s} - \delta, \hat{x}_\text{s} + \delta]$. Across this interface, for a pseudo-steady state, we compute the losses to balance:

\begin{align*}
    & \rho_1^* u_1^* = \rho_2^* u_2^*                           \\
    & \rho_1^* {u_1^*}^2 + p_1^* = \rho_2^* {u_2^*}^2 + p_2^*   \\
    & h_1^* + \frac{{u_1^*}^2}{2} = h_2^* + \frac{{u_2^*}^2}{2}
\end{align*}

where $h^*$ is the specific enthalpy. This is the only term that allows for the exchange of information between the two sub-networks, and to the interface itself.
